% !TEX root = ../../presentation.tex
% =============================================================
%  Chapter — Transmission line: creepage distance of DC OHL insulators
%  \input from presentation.tex as transmission-line/creepage-distance/creepage-distance
%
%  Figures in ./fig (cd_*.pdf), vector, generated by
%  fig/creepage_figures.py (numpy + matplotlib only).
%  Reference case: 400 kV pole-to-earth DC voltage.
%  USCD values are indicative ranges from IEC TS 60815-4 and
%  CIGRE TB 518 --- not agreed AmWatt design constants.
%  ESDD/NSDD measurement after IEC TS 60815-1 and IEC 60507.
% =============================================================
\newcommand{\cdfig}[1]{transmission-line/creepage-distance/fig/#1}

\sectionslide{Creepage Distance}{Insulator dimensioning for DC overhead lines}

% ---- Abbreviations -------------------------------------------
\begin{frame}{Abbreviations}
  \begin{columns}[T,onlytextwidth]
    \begin{column}{0.48\textwidth}
      {\small
      \begin{tabularx}{\linewidth}{@{}l@{\hspace{0.8em}}X@{}}
        AC    & Alternating current \\
        CIGRE & International Council on Large Electric Systems \\
        DC    & Direct current \\
        ESDD  & Equivalent salt deposit density \\
        HVDC  & High-voltage direct current \\
        IEC   & International Electrotechnical Commission \\
      \end{tabularx}}
    \end{column}
    \begin{column}{0.48\textwidth}
      {\small
      \begin{tabularx}{\linewidth}{@{}l@{\hspace{0.8em}}X@{}}
        NSDD  & Non-soluble deposit density \\
        OHL   & Overhead line \\
        SPS   & Site pollution severity \\
        TB    & Technical brochure \\
        TS    & Technical specification \\
        USCD  & Unified specific creepage distance \\
      \end{tabularx}}
    \end{column}
  \end{columns}
\end{frame}

% ---- Definition ----------------------------------------------
\begin{frame}{Creepage distance\\Definition and reference voltage}
  \begin{columns}[T,onlytextwidth]
    \begin{column}{0.56\textwidth}
      \begin{block}{Creepage distance}
        Shortest distance along the insulator surface between the two
        conductive parts that normally carry the operating voltage.
      \end{block}
      {\small
      \begin{itemize}
        \setlength{\itemsep}{0.4em}
        \item Pollution sets it: a wetted pollution layer conducts and can
              flash over at normal operating voltage.
        \item Specified as a USCD in mm/kV of the pole-to-earth DC
              operating voltage~[1].
        \item No single value: the USCD follows from site pollution and
              insulator type~[1], [2].
      \end{itemize}}
    \end{column}
    \begin{column}{0.40\textwidth}\centering
      \vspace{1.2em}
      \[ L = \mathrm{USCD} \cdot U_{dpe} \cdot C_{a} \]
      \vspace{0.4em}
      {\color{atGray}\footnotesize
      \begin{tabularx}{\linewidth}{@{}l@{\hspace{0.6em}}>{\raggedright\arraybackslash}X@{}}
        $L$         & total creepage in mm \\
        USCD        & unified specific creepage distance in mm/kV \\
        $U_{dpe}$   & pole-to-earth DC operating voltage in kV \\
        $C_{a}$     & altitude correction ($=1$ up to 1000\,m, see altitude slide) \\
      \end{tabularx}}
    \end{column}
  \end{columns}
\end{frame}

% ---- Why DC needs more creepage ------------------------------
\begin{frame}{Creepage distance\\Why DC needs more than AC}
  \begin{columns}[T,onlytextwidth]
    \begin{column}{0.42\textwidth}
      {\small
      \begin{itemize}
        \setlength{\itemsep}{0.4em}
        \item The steady DC field attracts charged particles, so a DC
              insulator collects roughly 1.5--2 times the pollution of an
              AC insulator at the same site~[2].
        \item A dry-band arc on DC has no current zero to extinguish it.
        \item For the same site, DC USCD values are therefore well above
              AC values.
      \end{itemize}}
    \end{column}
    \begin{column}{0.54\textwidth}
      {\small
      \renewcommand{\arraystretch}{1.25}
      \begin{tabularx}{\linewidth}{@{}>{\raggedright\arraybackslash}X
                                     >{\raggedright\arraybackslash}p{0.27\linewidth}
                                     >{\raggedright\arraybackslash}p{0.30\linewidth}@{}}
        \toprule
        \textbf{Aspect} & \multicolumn{1}{c}{\textbf{AC}} & \multicolumn{1}{c}{\textbf{DC}} \\
        \midrule
        Pollution build-up      & wind, gravity    & plus electrostatic attraction \\
        \atrowfill
        Dry-band arc            & current zero every half cycle & no current zero \\
        Ionic migration in glass and porcelain & negligible & needs DC-type glass, zinc sleeve on pin \\
        \atrowfill
        Relative USCD           & 1                & about 1.5--2 \\
        \bottomrule
      \end{tabularx}}
    \end{column}
  \end{columns}
\end{frame}

% ---- Site pollution: ESDD and NSDD ---------------------------
\begin{frame}{Site pollution\\ESDD and NSDD}
  \begin{columns}[T,onlytextwidth]
    \begin{column}{0.50\textwidth}
      {\small
      \begin{itemize}
        \setlength{\itemsep}{0.35em}
        \item \textbf{ESDD}: mass of NaCl per insulator surface area that,
              dissolved in water, gives the same conductivity as the
              soluble deposit washed off the insulator, in mg/cm$^2$~[3].
              The soluble deposit is the salt (sea, industrial, desert dust)
              that dissolves into ions when wetted and makes the layer conduct.
        \item \textbf{NSDD}: mass of the non-soluble residue (dust, sand,
              soil) per surface area, in mg/cm$^2$~[3]. It holds water and
              so helps the soluble salt to form a conducting layer.
        \item Both are measured by washing a reference insulator with
              demineralised water after a year or more on site.
      \end{itemize}}
    \end{column}
    \begin{column}{0.46\textwidth}\centering
      \setlength{\abovedisplayskip}{2pt}\setlength{\belowdisplayskip}{2pt}
      \[ \sigma_{20} = \sigma_{\theta} \cdot \left[1 - b \cdot (\theta - 20\,^{\circ}\mathrm{C})\right] \]
      \[ S_{a} = \left(5.7 \cdot \sigma_{20}\right)^{1.03} \]
      \[ \mathrm{ESDD} = \frac{S_{a} \cdot V}{A} \]
      \[ \mathrm{NSDD} = \frac{1000 \cdot (W_{f} - W_{i})}{A} \]
      \vspace{0.3em}
      {\color{atGray}\footnotesize
      \begin{tabularx}{\linewidth}{@{}l@{\hspace{0.6em}}>{\raggedright\arraybackslash}X@{}}
        $\sigma_{\theta}$, $\sigma_{20}$ & wash-water conductivity at $\theta$ and at 20\,$^{\circ}$C in S/m \\
        $b$         & temperature factor~[4] \\
        $S_{a}$     & salinity in kg/m$^3$ (= mg/cm$^3$) \\
        $V$, $A$    & wash-water volume in cm$^3$, washed area in cm$^2$ \\
        $W_{i}$, $W_{f}$ & filter mass before and after, dried, in g \\
      \end{tabularx}}
    \end{column}
  \end{columns}
\end{frame}

% ---- Site pollution: from ESDD to creepage --------------------
\begin{frame}{Site pollution\\From ESDD and NSDD to creepage}
  \begin{columns}[T,onlytextwidth]
    \begin{column}{0.52\textwidth}
      {\small
      \begin{itemize}
        \setlength{\itemsep}{0.35em}
        \item ESDD and NSDD together give the SPS class of the site,
              from very light to very heavy~[3].
        \item For DC the SPS must be measured on DC-energised insulators,
              or estimated from AC or non-energised data with a DC/AC
              accumulation ratio of about 1.5--2~[1], [2].
        \item Flashover strength falls with ESDD and NSDD as a power law,
              so the USCD rises with ESDD$^{a}$ and NSDD$^{b_{n}}$; typically
              $a \approx 0.2$--0.4, $b_{n} \approx 0.1$--0.2.
      \end{itemize}}
    \end{column}
    \begin{column}{0.44\textwidth}\centering
      \setlength{\abovedisplayskip}{3pt}\setlength{\belowdisplayskip}{3pt}
      \[ E_{f} = \frac{U_{f}}{L} = A_{f} \cdot \mathrm{ESDD}^{-a} \cdot \mathrm{NSDD}^{-b_{n}} \]
      \[ \begin{aligned}
        \mathrm{USCD}_{2} = \mathrm{USCD}_{1}
          &\cdot \left(\frac{\mathrm{ESDD}_{2}}{\mathrm{ESDD}_{1}}\right)^{a}\\
          &\cdot \left(\frac{\mathrm{NSDD}_{2}}{\mathrm{NSDD}_{1}}\right)^{b_{n}}
      \end{aligned} \]
      {\color{atGray}\footnotesize
      \begin{tabularx}{\linewidth}{@{}l@{\hspace{0.6em}}>{\raggedright\arraybackslash}X@{}}
        $E_{f}$     & flashover strength per creepage in kV/m \\
        $A_{f}$     & constant of the insulator type \\
        $a$, $b_{n}$ & ESDD and NSDD exponents from the supplier's DC
                      pollution tests or published data~[2] \\
      \end{tabularx}}\\[0.3em]
      {\footnotesize Example: doubling ESDD with $a = 0.3$ needs
      $2^{0.3} \approx 1.23$, i.e.\ 23\,\% more creepage.}
    \end{column}
  \end{columns}
\end{frame}

% ---- Typical USCD values -------------------------------------
\begin{frame}{Creepage distance\\Indicative USCD for DC OHL insulators}
  \begin{columns}[T,onlytextwidth]
    \begin{column}{0.54\textwidth}
      {\small
      \renewcommand{\arraystretch}{1.25}
      \begin{tabularx}{\linewidth}{@{}>{\raggedright\arraybackslash}X r r@{}}
        \toprule
        \textbf{Site pollution} & \multicolumn{1}{c}{\textbf{Ceramic / glass}}
                                & \multicolumn{1}{c}{\textbf{Composite}} \\
        \midrule
        Light               & 30--40\,mm/kV & 25--30\,mm/kV \\
        \atrowfill
        Medium              & 40--50\,mm/kV & 30--38\,mm/kV \\
        Heavy               & 50--60\,mm/kV & 38--45\,mm/kV \\
        \atrowfill
        Very heavy          & mostly too long & 45--55\,mm/kV \\
        \bottomrule
      \end{tabularx}}
    \end{column}
    \begin{column}{0.42\textwidth}
      {\small
      \begin{itemize}
        \setlength{\itemsep}{0.4em}
        \item Composite (silicone) insulators keep a water-repellent surface,
              so they need roughly 20--30\,\% less creepage.
        \item Most recent HVDC lines therefore use composite long-rod
              insulators.
        \item The ranges are a starting point; the design value comes from
              the measured site pollution severity.
      \end{itemize}}
    \end{column}
  \end{columns}
  \atnote{Indicative values from IEC TS 60815-4~[1] and CIGRE TB 518~[2];
  not agreed AmWatt design constants.}
\end{frame}

% ---- Worked example 400 kV -----------------------------------
\begin{frame}{Creepage distance\\Worked example at 400\,kV}
  \begin{columns}[T,onlytextwidth]
    \begin{column}{0.46\textwidth}
      {\small Study voltage $U_{dpe}=400\,\mathrm{kV}$, site below 1000\,m
      ($C_{a}=1$):}
      \[ \begin{aligned}
        L_{cer} &= 45\,\mathrm{mm/kV} \cdot 400\,\mathrm{kV} = 18.0\,\mathrm{m}\\[0.15em]
        L_{comp} &= 34\,\mathrm{mm/kV} \cdot 400\,\mathrm{kV} = 13.6\,\mathrm{m}
      \end{aligned} \]
      {\small
      \begin{itemize}
        \setlength{\itemsep}{0.3em}
        \item Medium pollution, mid-range USCD for each insulator type.
        \item With about 600\,mm creepage per DC fog-type disc, the ceramic
              string needs about 30 discs.
        \item Over the full light-to-heavy range: 12.0--24.0\,m (ceramic).
      \end{itemize}}
    \end{column}
    \begin{column}{0.50\textwidth}\centering
      \topoimg[width=\linewidth,height=0.62\textheight,keepaspectratio]{\cdfig{cd_creepage_voltage.pdf}}\\[0.2em]
      {\color{atGray}\footnotesize Total creepage versus pole voltage,
      medium-pollution bands.}
    \end{column}
  \end{columns}
\end{frame}

% ---- Altitude correction -------------------------------------
\begin{frame}{Creepage distance\\Altitude correction}
  \begin{columns}[T,onlytextwidth]
    \begin{column}{0.48\textwidth}
      {\small
      \begin{itemize}
        \setlength{\itemsep}{0.4em}
        \item Lower air pressure lowers the pollution flashover voltage, so
              above 1000\,m more creepage is needed.
        \item $m$ comes from pollution tests at reduced pressure and depends
              on insulator type, profile and pollution level; DC tests give
              about 0.3--0.5~[2].
        \item Example: 400\,kV ceramic at 2000\,m, $m=0.5$:
              $18.0\,\mathrm{m} \cdot 1.063 \approx 19.1\,\mathrm{m}$.
      \end{itemize}}
    \end{column}
    \begin{column}{0.48\textwidth}\centering
      \setlength{\abovedisplayskip}{3pt}\setlength{\belowdisplayskip}{3pt}
      \[ U_{f} = U_{f0} \cdot \left(\frac{p}{p_{0}}\right)^{m}
         \qquad \frac{p}{p_{0}} = e^{-H/8150\,\mathrm{m}} \]
      \[ C_{a} = e^{\,m \cdot \frac{H-1000\,\mathrm{m}}{8150\,\mathrm{m}}} \]
      \vspace{0.3em}
      {\small
      \renewcommand{\arraystretch}{1.2}
      \begin{tabular}{@{}r r r@{}}
        \toprule
        \multicolumn{1}{c}{\textbf{$H$}} & \multicolumn{1}{c}{\textbf{$m=0.35$}}
                                       & \multicolumn{1}{c}{\textbf{$m=0.5$}} \\
        \midrule
        1500\,m & 1.022 & 1.031 \\
        \atrowfill
        2000\,m & 1.044 & 1.063 \\
        2500\,m & 1.067 & 1.096 \\
        \bottomrule
      \end{tabular}}
    \end{column}
  \end{columns}
  \atnote{$U_{f}$, $U_{f0}$ pollution flashover voltage at pressure $p$ and at
  $p_{0}$ (sea level); $H$ altitude; $m$ empirical pressure exponent.}
\end{frame}

% ---- Design procedure ----------------------------------------
\begin{frame}{Creepage distance\\Design procedure}
  \begin{block}{From site to string length}
    Measure the pollution, select the USCD, correct it, then check the
    string against the overvoltage requirements.
  \end{block}
  {\small
  \begin{enumerate}
    \setlength{\itemsep}{0.4em}
    \item Determine the SPS along the route: ESDD and NSDD from test
          insulators exposed under DC voltage, or from service experience~[1], [2].
    \item Convert the SPS into a USCD with the IEC TS 60815-4 / CIGRE curves,
          or with the supplier's DC pollution test data for the chosen
          insulator profile.
    \item Apply the altitude correction $C_{a}$ above 1000\,m, and the
          corrections for insulator diameter and profile.
    \item Check the resulting string against the switching and lightning
          impulse requirements; the longer of the two governs.
  \end{enumerate}}
\end{frame}

% ---- References ----------------------------------------------
\begin{frame}{References}
  \noindent{\small
  \begin{enumerate}[{[}1{]}]
    \setlength{\itemsep}{0.5em}
    \item IEC TS 60815-4:2016, \emph{Selection and dimensioning of
          high-voltage insulators intended for use in polluted conditions ---
          Part 4: Insulators for d.c.\ systems}.
    \item CIGRE Working Group C4.303, TB 518, \emph{Outdoor insulation in
          polluted conditions: Guidelines for selection and dimensioning ---
          Part 2: The DC case}, 2012.
    \item IEC TS 60815-1:2008, \emph{Selection and dimensioning of
          high-voltage insulators intended for use in polluted conditions ---
          Part 1: Definitions, information and general principles}.
    \item IEC 60507:2013, \emph{Artificial pollution tests on high-voltage
          ceramic and glass insulators to be used on a.c.\ systems}.
  \end{enumerate}}
\end{frame}
